\section{Explicit Three-Qubit Interaction Coordinates}
\label{app:explicit_3q_invariants}

\noindent
For three qubits, the interaction coordinates can be evaluated directly from the diagonal phase values $\{\phi_{abc}\}_{a,b,c\in\{0,1\}}$ using the parity form of Eq.~\eqref{eq:phase_coefficients_parity_sum}. Writing $\vec{x}=(a,b,c)$, this yields
\begin{equation}
\label{eq:phiS_3q_general}
\phi(S)
=
\frac{1}{8}
\sum_{a,b,c\in\{0,1\}}
(-1)^{\sum_{i\in S} x_i}\,\phi_{abc},
\end{equation}
where $\phi_{abc} \equiv \phi(a,b,c)$ denotes the phase acquired by the computational-basis state $\ket{abc}$.

\medskip
\noindent
That is, each interaction coordinate is obtained as a parity-weighted sum over the measured phase values, where the sign for each computational-basis state is determined by the parity of the bits indexed by $S$. We list the explicit expressions for all nontrivial subsets below.

\medskip
\noindent\textbf{Three-body interaction coordinate $\phi(\{a,b,c\})$.}
\begin{equation}
\boxed{
\begin{aligned}
\phi(\{a,b,c\})
&= \frac{1}{8}\Big(
\phi_{000}
-\phi_{001}
-\phi_{010}
+\phi_{011} \\
&\qquad\quad
-\phi_{100}
+\phi_{101}
+\phi_{110}
-\phi_{111}
\Big)
\end{aligned}
}
\end{equation}

\medskip
\noindent\textbf{Two-body interaction coordinates.}
\begin{equation}
\boxed{
\begin{aligned}
\phi(\{a,b\})
&= \frac{1}{8}\Big(
\phi_{000}
+\phi_{001}
-\phi_{010}
-\phi_{011} \\
&\qquad\quad
-\phi_{100}
-\phi_{101}
+\phi_{110}
+\phi_{111}
\Big)
\end{aligned}
}
\end{equation}

\begin{equation}
\boxed{
\begin{aligned}
\phi(\{a,c\})
&= \frac{1}{8}\Big(
\phi_{000}
-\phi_{001}
+\phi_{010}
-\phi_{011} \\
&\qquad\quad
-\phi_{100}
+\phi_{101}
-\phi_{110}
+\phi_{111}
\Big)
\end{aligned}
}
\end{equation}

\begin{equation}
\boxed{
\begin{aligned}
\phi(\{b,c\})
&= \frac{1}{8}\Big(
\phi_{000}
-\phi_{001}
-\phi_{010}
+\phi_{011} \\
&\qquad\quad
+\phi_{100}
-\phi_{101}
-\phi_{110}
+\phi_{111}
\Big)
\end{aligned}
}
\end{equation}

\medskip
\noindent\textbf{Single-qubit interaction coordinates.}
\begin{equation}
\boxed{
\begin{aligned}
\phi(\{a\})
&= \frac{1}{8}\Big(
\phi_{000}
+\phi_{001}
+\phi_{010}
+\phi_{011} \\
&\qquad\quad
-\phi_{100}
-\phi_{101}
-\phi_{110}
-\phi_{111}
\Big)
\end{aligned}
}
\end{equation}

\begin{equation}
\boxed{
\begin{aligned}
\phi(\{b\})
&= \frac{1}{8}\Big(
\phi_{000}
+\phi_{001}
-\phi_{010}
-\phi_{011} \\
&\qquad\quad
+\phi_{100}
+\phi_{101}
-\phi_{110}
-\phi_{111}
\Big)
\end{aligned}
}
\end{equation}

\begin{equation}
\boxed{
\begin{aligned}
\phi(\{c\})
&= \frac{1}{8}\Big(
\phi_{000}
-\phi_{001}
+\phi_{010}
-\phi_{011} \\
&\qquad\quad
+\phi_{100}
-\phi_{101}
+\phi_{110}
-\phi_{111}
\Big)
\end{aligned}
}
\end{equation}


\section{Explicit Three-Qubit Reconstruction of the Phase Map}
\label{app:explicit_3q_phase_map_reconstruction}

\noindent
The diagonal representation introduced in the main text shows that the unitary $U_{\mathrm{diag}}$ is fully characterized by the phases $\phi(\vec{x})$ acquired by computational-basis states. Since global
phase is not observable, the phases $\phi(\vec{x})$ cannot be measured directly. Instead, experiments provide access to relative phases between computational-basis states, which can be obtained using conditional
Ramsey interferometry \cite{ramsey1950molecular, childress2006coherent, balasubramanian2009ultralong, taminiau2014universal}.

\medskip
\noindent
To measure such phase differences, one qubit $p$ is used as a probe while the remaining qubits are prepared in a fixed computational-basis configuration
\[
\ket{s}=\ket{x_1\dots x_{p-1}x_{p+1}\dots x_n}.
\]
The probe qubit is initialized in the superposition
\[
\ket{+}_p=\frac{\ket{0}_p+\ket{1}_p}{\sqrt2},
\]
yielding the joint state
\[
\frac{1}{\sqrt2}\bigl(\ket{0,s}+\ket{1,s}\bigr).
\]

\medskip
\noindent
After evolution under $U_{\mathrm{diag}}$, the two components accumulate different phases,
\[
\frac{1}{\sqrt2}
\left(
e^{i\phi(0,s)}\ket{0,s}
+
e^{i\phi(1,s)}\ket{1,s}
\right).
\]
Ramsey interferometry therefore measures the relative phase
\begin{equation}
\Theta_p(s)=\phi(1,s)-\phi(0,s),
\label{eq:ramsey_phase}
\end{equation}
which corresponds to the phase difference between computational-basis states that differ only in qubit $p$.

\medskip
\noindent
For a three-qubit subsystem with computational-basis phases
\[
\{\phi_{abc}\}_{a,b,c\in\{0,1\}},
\qquad
\phi_{abc}\equiv \phi(a,b,c),
\]
the conditional Ramsey measurements determine relative phases along the edges of the three-dimensional Boolean hypercube of computational-basis states.

\medskip
\noindent
Choosing qubit $a$ as probe yields
\begin{align}
\Theta_a(0,0) &= \phi_{100}-\phi_{000}, \\
\Theta_a(0,1) &= \phi_{101}-\phi_{001}, \\
\Theta_a(1,0) &= \phi_{110}-\phi_{010}, \\
\Theta_a(1,1) &= \phi_{111}-\phi_{011}.
\end{align}

\noindent
Choosing qubit $b$ as probe gives
\begin{align}
\Theta_b(0,0) &= \phi_{010}-\phi_{000}, \\
\Theta_b(0,1) &= \phi_{011}-\phi_{001}, \\
\Theta_b(1,0) &= \phi_{110}-\phi_{100}, \\
\Theta_b(1,1) &= \phi_{111}-\phi_{101}.
\end{align}

\noindent
Choosing qubit $c$ as probe gives
\begin{align}
\Theta_c(0,0) &= \phi_{001}-\phi_{000}, \\
\Theta_c(0,1) &= \phi_{011}-\phi_{010}, \\
\Theta_c(1,0) &= \phi_{101}-\phi_{100}, \\
\Theta_c(1,1) &= \phi_{111}-\phi_{110}.
\end{align}

\medskip
\noindent
These twelve measurements provide linear constraints on the eight unknown phases $\phi_{abc}$. Since the phase map is defined only up to a global phase, we fix the reference
\[
\phi_{000}=0.
\]
The remaining phases can then be reconstructed by solving the resulting linear system. One convenient reconstruction is
\begin{align}
\phi_{100} &= \Theta_a(0,0), \\
\phi_{010} &= \Theta_b(0,0), \\
\phi_{001} &= \Theta_c(0,0), \\
\phi_{110} &= \phi_{010}+\Theta_a(1,0), \\
\phi_{101} &= \phi_{001}+\Theta_a(0,1), \\
\phi_{011} &= \phi_{001}+\Theta_b(0,1), \\
\phi_{111} &= \phi_{011}+\Theta_a(1,1).
\end{align}

\noindent
Equivalently, the final phase can be reconstructed along different paths on the hypercube, for example
\begin{equation}
\phi_{111}=\phi_{110}+\Theta_c(1,1)=\phi_{101}+\Theta_b(1,1),
\end{equation}
which is consistent provided the measured phase differences satisfy the
hypercube compatibility relations.

\medskip
\noindent
The same reconstruction principle extends directly to an $n$-qubit system. In that case the computational-basis phases $\{\phi(\vec{x})\}_{\vec{x}\in\{0,1\}^n}$ correspond to the vertices of an $n$-dimensional Boolean hypercube. Conditional Ramsey measurements determine phase differences along its edges,
\[
\Theta_p(s)=\phi(1,s)-\phi(0,s),
\]
for each probe qubit $p$ and spectator configuration $s$. Repeating the experiment for all $p$ and $s$ yields $n2^{\,n-1}$ linear relations among the $2^n$ phases. Fixing a reference phase (e.g.\ $\phi(0,\dots,0)=0$)
removes the global phase freedom, after which the remaining phases can be reconstructed by solving the resulting linear system.

\medskip
\noindent
Once the full set $\{\phi(\vec{x})\}$ has been reconstructed, the corresponding interaction coordinates follow from the Walsh-Hadamard parity sums introduced in the main text and listed explicitly for the three-qubit
case in Appendix~\ref{app:explicit_3q_invariants}.

\end{document}
